\documentclass[UTF8,a4paper,11pt,openany]{ctexbook}
\usepackage{../common/statlearnbook}
\SLSetBookMetadata{第 3 册}{统计学习方法讲义 v2.6.0}{优化模型与序列模型}
\begin{document}
\setcounter{page}{338}
\setcounter{chapter}{17}
\setcounter{figure}{0}
\pagestyle{headings}
\chaptermark{最大熵模型}
\begin{probabilitybox}
熵越大，概率质量越不集中；新增可靠约束后，可行集合缩小，最大熵解才相应偏离均匀分布。
\end{probabilitybox}
% v2.6.0 figure UID: FIG-P282-01
% Image-reference reverse redraw: preserve the simplex data and strengthen the feasible-point reading path.
\tikzset{
  slfig-FIG-P282-01/.style={
    every path/.append style={line cap=round,line join=round},
    font=\fontsize{9.2}{11}\selectfont,
  },
  slfig-FIG-P282-01-focus-label/.style={
    slfig direct label,font=\fontsize{9.4}{11.2}\selectfont,
    text=SLGold!82!black,fill=none,inner sep=0pt,
  },
  slfig-FIG-P282-01-contour-label/.style={
    font=\fontsize{9}{10.6}\selectfont,text=SLInk!70,
    fill=none,inner sep=0pt,
  },
  slfig-FIG-P282-01-constraint-label/.style={
    font=\fontsize{9.2}{11}\selectfont,fill=none,inner sep=0pt,
  },
  slfig-FIG-P282-01-vertex/.style={font=\fontsize{9.6}{11.4}\selectfont,text=SLInk!92},
}
\pgfplotsset{
  slfig-FIG-P282-01-axis/.style={clip=false},
  slfig-FIG-P282-01-simplex/.style={draw=SLBlue!92,line width=1.00pt,fill=SLBlue,fill opacity=.035},
  slfig-FIG-P282-01-contour/.style={contour/draw color=SLInk!64,line width=.56pt,smooth},
  slfig-FIG-P282-01-full-one/.style={draw=SLBlue!52,line width=.52pt,densely dashed},
  slfig-FIG-P282-01-full-two/.style={draw=SLTeal!58,line width=.52pt,dash dot},
  slfig-FIG-P282-01-feasible-one/.style={draw=SLBlue,line width=1.00pt},
  slfig-FIG-P282-01-feasible-two/.style={draw=SLTeal,line width=1.00pt,dash pattern=on 6pt off 2pt},
}
% Entropy contours below are numerical solutions of H(p)=-sum_i p_i ln p_i.
% Explicit barycentric transform: x=p_2+p_3/2, y=(sqrt(3)/2)p_3.
\begin{figure}[H]
\centering
\begin{tikzpicture}[slfig-FIG-P282-01]
\begin{axis}[slfig axis,slfig-FIG-P282-01-axis,
  width=.78\linewidth,height=.69\linewidth,
  xmin=-.12,xmax=1.38,ymin=-.12,ymax=.99,
  axis equal image,hide axis,clip=false,
]
  \addplot[slfig-FIG-P282-01-simplex]
    coordinates {(0,0) (.5,.8660254) (1,0)}\closedcycle;

  \addplot[
    contour prepared={labels=false},slfig-FIG-P282-01-contour,
  ] table[
    x expr=\thisrow{p2}+.5*\thisrow{p3},
    y expr=.8660254*\thisrow{p3},
    z=H,
  ] {
    p1 p2 p3 H
    % Numerically solved -sum p_i ln p_i = 0.80
    0.6221467 0.0445199 0.3333333 0.80
    0.6833946 0.0770707 0.2395347 0.80
    0.7119860 0.1440070 0.1440070 0.80
    0.6833946 0.2395347 0.0770707 0.80
    0.6221467 0.3333333 0.0445199 0.80
    0.5545696 0.4143114 0.0311190 0.80
    0.4862547 0.4862547 0.0274905 0.80
    0.4143114 0.5545696 0.0311190 0.80
    0.3333333 0.6221467 0.0445199 0.80
    0.2395347 0.6833946 0.0770707 0.80
    0.1440070 0.7119860 0.1440070 0.80
    0.0770707 0.6833946 0.2395347 0.80
    0.0445199 0.6221467 0.3333333 0.80
    0.0311190 0.5545696 0.4143114 0.80
    0.0274905 0.4862547 0.4862547 0.80
    0.0311190 0.4143114 0.5545696 0.80
    0.0445199 0.3333333 0.6221467 0.80
    0.0770707 0.2395347 0.6833946 0.80
    0.1440070 0.1440070 0.7119860 0.80
    0.2395347 0.0770707 0.6833946 0.80
    0.3333333 0.0445199 0.6221467 0.80
    0.4143114 0.0311190 0.5545696 0.80
    0.4862547 0.0274905 0.4862547 0.80
    0.5545696 0.0311190 0.4143114 0.80
    0.6221467 0.0445199 0.3333333 0.80

    % Numerically solved -sum p_i ln p_i = 0.95
    0.5471779 0.1194888 0.3333333 0.95
    0.5846883 0.1493287 0.2659830 0.95
    0.6002461 0.1998769 0.1998769 0.95
    0.5846883 0.2659830 0.1493287 0.95
    0.5471779 0.3333333 0.1194888 0.95
    0.5004902 0.3945170 0.1049929 0.95
    0.4496385 0.4496385 0.1007230 0.95
    0.3945170 0.5004902 0.1049929 0.95
    0.3333333 0.5471779 0.1194888 0.95
    0.2659830 0.5846883 0.1493287 0.95
    0.1998769 0.6002461 0.1998769 0.95
    0.1493287 0.5846883 0.2659830 0.95
    0.1194888 0.5471779 0.3333333 0.95
    0.1049929 0.5004902 0.3945170 0.95
    0.1007230 0.4496385 0.4496385 0.95
    0.1049929 0.3945170 0.5004902 0.95
    0.1194888 0.3333333 0.5471779 0.95
    0.1493287 0.2659830 0.5846883 0.95
    0.1998769 0.1998769 0.6002461 0.95
    0.2659830 0.1493287 0.5846883 0.95
    0.3333333 0.1194888 0.5471779 0.95
    0.3945170 0.1049929 0.5004902 0.95
    0.4496385 0.1007230 0.4496385 0.95
    0.5004902 0.1049929 0.3945170 0.95
    0.5471779 0.1194888 0.3333333 0.95

    % Numerically solved -sum p_i ln p_i = 1.05
    0.4590562 0.2076105 0.3333333 1.05
    0.4775120 0.2277872 0.2947008 1.05
    0.4844746 0.2577627 0.2577627 1.05
    0.4775120 0.2947008 0.2277872 1.05
    0.4590562 0.3333333 0.2076105 1.05
    0.4334471 0.3699775 0.1965754 1.05
    0.4034414 0.4034414 0.1931172 1.05
    0.3699775 0.4334471 0.1965754 1.05
    0.3333333 0.4590562 0.2076105 1.05
    0.2947008 0.4775120 0.2277872 1.05
    0.2577627 0.4844746 0.2577627 1.05
    0.2277872 0.4775120 0.2947008 1.05
    0.2076105 0.4590562 0.3333333 1.05
    0.1965754 0.4334471 0.3699775 1.05
    0.1931172 0.4034414 0.4034414 1.05
    0.1965754 0.3699775 0.4334471 1.05
    0.2076105 0.3333333 0.4590562 1.05
    0.2277872 0.2947008 0.4775120 1.05
    0.2577627 0.2577627 0.4844746 1.05
    0.2947008 0.2277872 0.4775120 1.05
    0.3333333 0.2076105 0.4590562 1.05
    0.3699775 0.1965754 0.4334471 1.05
    0.4034414 0.1931172 0.4034414 1.05
    0.4334471 0.1965754 0.3699775 1.05
    0.4590562 0.2076105 0.3333333 1.05
  };

  % Full constraint lines (light) and their simplex-feasible segments (dark).
  \addplot[slfig-FIG-P282-01-full-one]
    coordinates {(.638,-.1005) (.232,.6028)};
  \addplot[slfig-FIG-P282-01-full-two]
    coordinates {(.3077,-.0774) (.5757,.8510)};
  \addplot[slfig-FIG-P282-01-feasible-one]
    coordinates {(.58,0) (.29,.5023)};
  \addplot[slfig-FIG-P282-01-feasible-two]
    coordinates {(.33,0) (.5533,.7736)};

  \addplot[only marks,mark=*,mark size=3.1pt,draw=SLGold!82!black,fill=SLGold]
    coordinates {(.4133333,.2886751)};
  \draw[draw=SLGold!82!black,line width=.55pt]
    (axis cs:.4133333,.2886751)--(axis cs:.4133333,.08)--(axis cs:1.00,.08)--(axis cs:1.00,.18);
  \node[slfig-FIG-P282-01-focus-label,anchor=west,align=left]
    at (axis cs:1.04,.20) {唯一可行\\最大熵点};
  \node[slfig-FIG-P282-01-contour-label,anchor=west,align=left]
    at (axis cs:1.04,.62) {等高线\\外层：$H=0.80$\\中层：$H=0.95$\\内层：$H=1.05$};
  \node[slfig-FIG-P282-01-constraint-label,text=SLBlue,anchor=east]
    at (axis cs:.23,.64) {约束 1};
  \node[slfig-FIG-P282-01-constraint-label,text=SLTeal!82!black,anchor=west]
    at (axis cs:.67,.82) {约束 2};
  \node[slfig-FIG-P282-01-vertex,anchor=north east] at (axis cs:-.01,-.02) {$p_1=1$};
  \node[slfig-FIG-P282-01-vertex,anchor=north west] at (axis cs:1.01,-.02) {$p_2=1$};
  \node[slfig-FIG-P282-01-vertex,anchor=south] at (axis cs:.5,.88) {$p_3=1$};
\end{axis}
\end{tikzpicture}
\caption{概率单纯形上的熵等高线与线性约束；本例约束交点是唯一可行的最大熵点}\label{fig:V3-C01-simplex}
\end{figure}

图\ref{fig:V3-C01-simplex}把概率归一化表示为单纯形，把已知统计量表示为线性约束；可行集是两者的交集，最大熵解只在该交集内选取。
\paragraph{首次阅读检查：最大熵模型的定义。}
先识别对象类型与定义域，再逐项核对量词、支持集和特殊约定；例子只用于检验，不能替代定义。
\end{document}
